\chapter{Introduction}

The valuation of options is one of the central problems in quantitative financial mathematics. While closed-form solutions exist for simple European options -- most notably within the Black--Scholes framework -- the pricing of more complex derivatives is considerably more challenging. In particular, so-called basket options, whose payoff depends on multiple underlyings, generally require numerical methods such as Monte Carlo simulations. As the number of underlying assets increases, the computational effort grows substantially, since correlations between the assets must also be taken into account.

In recent years, machine learning methods have emerged as a powerful alternative to classical analytical and numerical valuation techniques. Instead of explicitly specifying a stochastic model for the dynamics of the underlyings, neural networks learn the functional relationship between input variables and the option price directly. After the training phase, the valuation of new configurations can be performed with low computational cost. Studies such as \cite{culkin2017machine} and \cite{ferguson2018deeply} demonstrate that deep neural networks are capable of approximating even complex basket option prices with high accuracy.

Parallel to these developments, quantum computing has emerged as a promising field of research. Variational quantum circuits (VQCs) represent a hybrid approach in which a parameterized quantum circuit is trained using classical optimization procedures. Due to their structurally different functional representation, VQCs differ fundamentally from classical neural networks. While classical models realize nonlinearity through explicit activation functions and typically produce piecewise linear approximations, nonlinear dependencies in a VQC arise implicitly from the parameterized unitary transformation and the subsequent measurement process.

In the context of basket options, the question arises whether the inherent ability of entangled quantum states to represent non-separable dependencies can provide advantages in modeling correlated financial quantities. Since the payoff of a basket option depends crucially on the joint distribution of multiple underlyings, it is natural to investigate to what extent the global interactions between qubits enabled by entanglement can contribute to capturing such dependency structures.

The objective of this thesis is to investigate the approximation of basket option prices using a VQC. As freely available historical datasets of actually traded basket options are only very limited, data generation is based on a simulation model. Volatilities and correlations are estimated from real historical time series and used within a Monte Carlo simulation framework. The resulting prices therefore represent model-based target values. The datasets generated in this way serve as the training and test basis for the models under consideration.

The focus lies on the comparison between a quantum machine learning model in the form of a VQC and established classical neural network architectures following \cite{culkin2017machine} and \cite{ferguson2018deeply}. Different dataset sizes as well as various basket payoff structures (Worst-of, Best-of, Average) are considered, and the influence of noise is also examined.

The thesis is structured as follows:
First, the theoretical foundations of basket options, Monte Carlo simulation, classical neural networks, and VQCs are presented. Subsequently, the methodology for data generation, model architecture, and training procedure is described. Building upon this, the experimental results are presented and discussed with regard to the performance of the quantum- and classical-based models. Finally, the findings are summarized and possible directions for future research are outlined.